\chapter{自动机}
\label{chap:automata}

自动机（automaton）用有限状态描述对输入串的逐位计算。本章先研究确定有限
自动机，再引入非确定有限自动机。

\section{确定有限自动机}

\begin{example}[奇偶校验]
  \label{ex:parity-dfa}
  对 $x=x_0x_1\cdots x_{n-1}\in\bitsstar$，定义

  \[
    \operatorname{PARITY}(x)=x_0\mathbin{\mathsf{xor}}x_1
      \mathbin{\mathsf{xor}}\cdots\mathbin{\mathsf{xor}}x_{n-1}.
  \]

  并约定 $\operatorname{PARITY}(\eps)=0$。

  下面的自动机中，$q_0$ 表示目前读到偶数个 $1$，$q_1$ 表示目前读到
  奇数个 $1$。接受状态为 $q_1$。

  \begin{center}
    \begin{tikzpicture}[
      shorten >=1pt,
      node distance=3.2cm,
      on grid,
      auto,
      initial text={}
    ]
      \node[state,initial] (q0) {$q_0$};
      \node[state,accepting] (q1) [right=of q0] {$q_1$};
      \path[->]
        (q0) edge[loop above] node {$0$} ()
             edge[bend left] node {$1$} (q1)
        (q1) edge[loop above] node {$0$} ()
             edge[bend left] node {$1$} (q0);
    \end{tikzpicture}
  \end{center}
\end{example}

\begin{definition}[确定有限自动机]
  \label{def:dfa}
  一个\emph{确定有限自动机}（deterministic finite automaton, DFA）是四元组

  \[
    M=(K,s,F,\delta),
  \]

  其中：

  \begin{enumerate}[label=\arabic*.]
    \item $K$ 是有限的状态集合（state set）；
    \item $s\in K$ 是初始状态（initial state）；
    \item $F\subseteq K$ 是接受状态（accepting state）的集合；
    \item $\delta:K\times\bits\to K$ 是转移函数（transition function）。
  \end{enumerate}

  对输入 $x=x_0x_1\cdots x_{n-1}$，自动机产生状态序列

  \[
    s_0=s,
    \qquad
    s_{i+1}=\delta(s_i,x_i)\quad(0\leq i<n).
  \]

  若 $s_n\in F$，则称 $M$ 接受 $x$；否则称 $M$ 拒绝 $x$。
\end{definition}

\begin{definition}[计算函数与判定语言]
  \label{def:dfa-compute-decide}
  设 $f:\bitsstar\to\bits$，$A\subseteq\bitsstar$。

  \begin{itemize}
    \item 如果对每个 $x$，$M$ 接受 $x$ 当且仅当 $f(x)=1$，则称
      $M$ \emph{计算} $f$；
    \item 如果对每个 $x$，$M$ 接受 $x$ 当且仅当 $x\in A$，则称
      $M$ \emph{判定} $A$。
  \end{itemize}
\end{definition}

\begin{definition}[DFA 的语言与正则语言]
  \label{def:regular-language}
  DFA $M$ 接受的全部字符串构成它的语言

  \[
    L(M)=\{x\in\bitsstar:M\text{ 接受 }x\}.
  \]

  如果存在 DFA $M$ 使得 $A=L(M)$，则称语言 $A$ 是\emph{正则的}
  （regular）。相应地，若 $A_f$ 正则，则称布尔函数 $f$ 是正则的。
\end{definition}

\begin{proposition}
  \label{prop:nonregular-language-exists}
  存在非正则语言。
\end{proposition}

\begin{proof}
  每个 DFA 都能编码为有限二进制串，所以全部 DFA 至多可数。
  另一方面，全部语言不可数：若能把它们枚举为
  $A_0,A_1,A_2,\ldots$，并把
  $\bitsstar$ 枚举为 $w_0,w_1,w_2,\ldots$，则语言

  \[
    D=\{w_i:w_i\notin A_i\}
  \]

  与每个 $A_i$ 都不同，矛盾。因此必有语言不能被任何 DFA 判定。
\end{proof}

\section{正则语言的封闭性}

\begin{theorem}[对并集封闭]
  \label{thm:regular-closed-under-union}
  若 $A$ 和 $B$ 都是正则语言，则 $A\cup B$ 也是正则语言。
\end{theorem}

\begin{proof}
  设

  \[
    M_A=(K_A,s_A,F_A,\delta_A),
    \qquad
    M_B=(K_B,s_B,F_B,\delta_B)
  \]

  分别判定 $A$ 和 $B$。构造乘积自动机 $M=(K,s,F,\delta)$，其中

  \begin{align*}
    K&=K_A\times K_B,\\
    s&=(s_A,s_B),\\
    F&=(F_A\times K_B)\cup(K_A\times F_B),\\
    \delta\bigl((q_A,q_B),a\bigr)
      &=\bigl(\delta_A(q_A,a),\delta_B(q_B,a)\bigr).
  \end{align*}

  归纳可知，$M$ 读完任意输入 $x$ 后的两个状态分量，恰好分别是
  $M_A$ 与 $M_B$ 读完 $x$ 后的状态。因此

  \[
    M\text{ 接受 }x
    \iff M_A\text{ 接受 }x\text{ 或 }M_B\text{ 接受 }x
    \iff x\in A\cup B.
  \]
\end{proof}

对语言 $A,B\subseteq\bitsstar$，定义它们的\emph{拼接}（concatenation）为

\[
  AB=\{xy:x\in A,\ y\in B\}.
\]

\begin{theorem}[对拼接封闭]
  \label{thm:regular-closed-under-concatenation}
  若 $A$ 和 $B$ 都是正则语言，则 $AB$ 也是正则语言。
\end{theorem}

直接用 DFA 构造拼接并不方便。完整证明将在后文给出；这里先引入 NFA。

\section{非确定有限自动机}

与 DFA 相比，NFA 可以从同一状态沿同一输入转移到多个状态，也可以进行不消耗
输入符号的 $\eps$-转移。

\begin{definition}[非确定有限自动机]
  \label{def:nfa}
  一个\emph{非确定有限自动机}（nondeterministic finite automaton, NFA）是
  四元组

  \[
    N=(K,s,F,\Delta),
  \]

  其中 $K,s,F$ 与 DFA 中的含义相同，而转移关系满足

  \[
    \Delta\subseteq K\times(\bits\cup\{\eps\})\times K.
  \]

  NFA 接受字符串 $x$，是指存在一条从 $s$ 到某个接受状态的路径，使得路径
  标签依次串联并删除所有 $\eps$ 后恰好得到 $x$。
\end{definition}

\begin{example}
  \label{ex:nfa-penultimate-one}
  语言

  \[
    A=\{w\in\bitsstar:\card{w}\geq2
      \text{ 且 }w\text{ 的倒数第二位为 }1\}
  \]

  可以由下面的 NFA 判定。在状态 $q_0$ 读到 $1$ 时，自动机既可以留在
  $q_0$，也可以猜测这个 $1$ 是倒数第二位并转移到 $q_1$；若随后恰好再读
  一个比特，便会到达接受状态。

  \begin{center}
    \begin{tikzpicture}[
      shorten >=1pt,
      node distance=2.8cm,
      on grid,
      auto,
      initial text={}
    ]
      \node[state,initial] (q0) {$q_0$};
      \node[state] (q1) [right=of q0] {$q_1$};
      \node[state,accepting] (q2) [right=of q1] {$q_2$};
      \path[->]
        (q0) edge[loop above] node {$0,1$} ()
             edge node {$1$} (q1)
        (q1) edge node {$0,1$} (q2);
    \end{tikzpicture}
  \end{center}

\end{example}

\lecture
\label{lec:nfa-and-regular-expressions}

\subsection{NFA 与 DFA 的等价性}

\begin{theorem}[NFA 与 DFA 等价]
  \label{thm:nfa-dfa-equivalence}
  对每个 NFA，都存在一个接受相同语言的 DFA；反之亦然。
\end{theorem}

\begin{proof}
  DFA 是每一步都只有唯一后继状态、且没有 $\eps$-转移的特殊 NFA，故一个
  DFA 可以直接看作 NFA。

  反过来，设 NFA 为 $N=(K,s,F,\Delta)$。对 $S\subseteq K$，记
  $E(S)$ 为从 $S$ 中某个状态出发，只经过有限次 $\eps$-转移所能到达的全部
  状态。构造 DFA

  \[
    D=(\mathcal{P}(K),E(\{s\}),F_D,\delta_D),
  \]

  其中 $\mathcal{P}(K)$ 是 $K$ 的幂集（power set），并令

  \begin{align*}
    F_D
      &=\{S\subseteq K:S\cap F\neq\varnothing\},\\
    \delta_D(S,a)
      &=E\bigl(\{q'\in K:\text{存在 }q\in S,
        (q,a,q')\in\Delta\}\bigr).
  \end{align*}

  这个构造称为状态子集构造（subset construction）。对输入长度归纳可知，
  $D$ 读完任意字符串 $x$ 后所处的状态，恰好是 $N$ 读完 $x$ 后所有可能状态
  的集合。因此该集合与 $F$ 相交，当且仅当 $N$ 有一条接受 $x$ 的计算路径。
  又因为 $K$ 有限，$\mathcal{P}(K)$ 也有限，所以 $D$ 确实是 DFA。
\end{proof}

\begin{corollary}
  \label{cor:regular-iff-nfa}
  一个语言是正则语言，当且仅当它能被某个 NFA 判定。
\end{corollary}

\begin{proof}
  由\cref{def:regular-language}和\cref{thm:nfa-dfa-equivalence}立即得到。
\end{proof}

\subsection{拼接与 Kleene 星}

\begin{proof}[\cref{thm:regular-closed-under-concatenation}的证明]
  取分别判定 $A$ 和 $B$ 的 NFA

  \[
    N_A=(K_A,s_A,F_A,\Delta_A),
    \qquad
    N_B=(K_B,s_B,F_B,\Delta_B),
  \]

  不妨设 $K_A\cap K_B=\varnothing$。构造 NFA

  \[
    N=\bigl(K_A\cup K_B,s_A,F_B,\Delta\bigr),
  \]

  其中

  \[
    \Delta=\Delta_A\cup\Delta_B
      \cup\{(q,\eps,s_B):q\in F_A\}.
  \]

  也就是说，从 $N_A$ 的每个接受状态增加一条到 $N_B$ 初始状态的
  $\eps$-转移。于是 $N$ 接受 $w$，当且仅当可以把 $w$ 写成 $w=xy$，使得
  $N_A$ 接受 $x$ 且 $N_B$ 接受 $y$。因此 $L(N)=AB$，再由
  \cref{cor:regular-iff-nfa}可知 $AB$ 正则。
\end{proof}

对语言 $A$，递归定义 $A^0=\{\eps\}$、$A^{n+1}=A^nA$，并称

\[
  A^*=\bigcup_{n\geq0}A^n
\]

为 $A$ 的 Kleene 星（Kleene star）。

\begin{theorem}[对 Kleene 星封闭]
  \label{thm:regular-closed-under-star}
  若 $A$ 是正则语言，则 $A^*$ 也是正则语言。
\end{theorem}

\begin{proof}
  取判定 $A$ 的 NFA $N_A=(K,s,F,\Delta)$，并增加一个新状态 $s'$。令

  \begin{align*}
    K'&=K\cup\{s'\},\\
    F'&=F\cup\{s'\},\\
    \Delta'&=\Delta\cup\{(s',\eps,s)\}
      \cup\{(q,\eps,s):q\in F\}.
  \end{align*}

  则 $N'=(K',s',F',\Delta')$ 可以先从 $s'$ 进入 $N_A$，并在每次到达
  $N_A$ 的接受状态后重新开始下一轮。新初态 $s'$ 本身是接受状态，所以
  $N'$ 也接受空串。因此 $L(N')=A^*$，由\cref{cor:regular-iff-nfa}可知
  $A^*$ 正则。
\end{proof}

\section{正则表达式}

\begin{definition}[正则表达式]
  \label{def:regular-expression}
  二进制字母表上的正则表达式（regular expression）按下面的规则递归定义：

  \begin{enumerate}[label=\arabic*.]
    \item $\varnothing$、$0$ 和 $1$ 是正则表达式，它们分别描述语言
      $\varnothing$、$\{0\}$ 和 $\{1\}$；
    \item 如果 $R_1,R_2$ 是正则表达式，那么 $R_1\cup R_2$ 和 $R_1R_2$
      也是正则表达式，并且
      \[
        L(R_1\cup R_2)=L(R_1)\cup L(R_2),
        \qquad
        L(R_1R_2)=L(R_1)L(R_2);
      \]
    \item 如果 $R$ 是正则表达式，那么 $R^*$ 也是正则表达式，并且
      \[
        L(R^*)=L(R)^*.
      \]
  \end{enumerate}

  在省略括号时，Kleene 星的优先级最高，拼接次之，并集最低。
  另外用 $\eps$ 作为 $\varnothing^*$ 的简写。
\end{definition}

\begin{example}
  \label{ex:regex-at-least-two-zeroes}
  至少含两个 $0$ 的二进制串构成的语言可以写成

  \[
    (0\cup1)^*0(0\cup1)^*0(0\cup1)^*.
  \]
\end{example}

\begin{theorem}[正则表达式刻画正则语言]
  \label{thm:regular-expression-characterization}
  一个语言是正则语言，当且仅当它能由正则表达式描述。
\end{theorem}

\begin{proof}
  先设语言由正则表达式 $R$ 描述。对 $R$ 的递归结构归纳：三个基本表达式
  描述的语言都是正则的；归纳步骤则分别由正则语言对并集、拼接和 Kleene 星
  的封闭性得到。因此 $L(R)$ 正则。

  反过来，设正则语言 $A$ 由某个 NFA $N$ 判定。先给 $N$ 增加一个没有入边
  的新初态和一个没有出边的唯一新接受态：从新初态向原初态添加
  $\eps$-转移，再从每个原接受态向新接受态添加 $\eps$-转移。

  把任意两个状态之间的所有边合并为一条边，并用这些边标签的并集作为新标签；
  没有边时，标签记为 $\varnothing$。随后逐个删除初态与接受态之外的状态。
  删除状态 $q$ 时，对每对保留的状态 $p,r$，将边 $p\to r$ 的标签更新为

  \[
    R_{pr}\ \cup\ R_{pq}(R_{qq})^*R_{qr}.
  \]

  该式分别涵盖了不经过 $q$ 的路径，以及进入 $q$、在 $q$ 处循环有限次后
  离开的路径。归纳可知，每次删除后，边标签仍描述原自动机中相应的全部路径。
  最后只剩新初态和新接受态，两者之间的边标签就是一个描述 $A$ 的正则表达式。
\end{proof}

\lecture
\label{lec:pumping-lemma}

\section{正则语言的泵引理}

\begin{theorem}[正则语言的泵引理（Pumping Theorem）]
  \label{thm:pumping-lemma-regular-languages}
  若 $L$ 是正则语言，则存在整数 $p\geq1$，使得每个满足
  $w\in L$ 且 $\card{w}\geq p$ 的字符串都可以写成 $w=xyz$，并满足：

  \begin{enumerate}[label=\arabic*.]
    \item 对每个整数 $i\geq0$，都有 $xy^iz\in L$；
    \item $\card{y}>0$；
    \item $\card{xy}\leq p$。
  \end{enumerate}

  这样的 $p$ 称为语言 $L$ 的一个泵长度（pumping length）。
\end{theorem}

\begin{proof}
  设 DFA $M=(K,s,F,\delta)$ 判定 $L$，取 $p=\card{K}$。令
  $w=w_0w_1\cdots w_{n-1}\in L$，其中 $n\geq p$；再令 $q_j$ 为 $M$
  读完 $w$ 的前 $j$ 位后所处的状态。

  序列 $q_0,q_1,\ldots,q_p$ 含有 $p+1$ 个状态。由抽屉原理，存在
  $0\leq r<t\leq p$ 使得 $q_r=q_t$。
  取

  \[
    x=w_0\cdots w_{r-1},\qquad
    y=w_r\cdots w_{t-1},\qquad
    z=w_t\cdots w_{n-1}.
  \]

  这里空的下标区间表示空串。

  因为 $r<t\leq p$，所以 $\card{y}=t-r>0$ 且 $\card{xy}=t\leq p$。
  又因为 $q_r=q_t$，从 $q_r$ 读入 $y$ 后仍回到 $q_r$。因此无论重复
  $y$ 多少次，读入 $z$ 之前自动机都处于 $q_r$；于是对每个 $i\geq0$，
  $M$ 都接受 $xy^iz$。
\end{proof}

\begin{example}
  \label{ex:pumping-length-simple-language}
  对正则语言

  \[
    L=01^*0,
  \]

  $p=3$ 是一个泵长度：若 $w=01^n0$ 且 $n\geq1$，取
  $x=0$、$y=1$、$z=1^{n-1}0$ 即可。另一方面，$p=2$ 不是泵长度，
  因为任取满足 $00=xyz$、$\card{y}>0$ 和 $\card{xy}\leq2$ 的分解，
  令 $i=0$ 都会得到长度小于 $2$ 的串，因而不属于 $L$。
\end{example}

\begin{example}[$\{0^n1^n:n\geq0\}$ 不是正则语言]
  \label{ex:zero-n-one-n-nonregular}
  令

  \[
    A=\{0^n1^n:n\geq0\}.
  \]

  假设 $A$ 正则，并令 $p$ 为\cref{thm:pumping-lemma-regular-languages}
  给出的泵长度。取 $w=0^p1^p\in A$。任取满足
  $w=xyz$、$\card{y}>0$ 且 $\card{xy}\leq p$ 的分解；由于 $xy$
  完全位于前 $p$ 个 $0$ 中，必有 $y=0^k$，其中 $1\leq k\leq p$。
  令 $i=0$，则

  \[
    xy^0z=0^{p-k}1^p\notin A,
  \]

  与泵引理矛盾。因此 $A$ 不是正则语言。
\end{example}
